\chapter{Tasarım Modeli}

\section{Genel Bakış}
Tasarım modeli, sistemin teknik mimarisini ve bileşenlerini detaylandırır. Bu model, mikroservis mimarisine dayalı sistemin yapısal tasarımını ve servisler arası etkileşimleri tanımlar.

\section{Component Diagram (Bileşen Diyagramı)}
Component diagram, sistemin yazılım bileşenlerini ve aralarındaki bağımlılıkları gösterir.

\begin{figure}[H]
    \centering
    \includegraphics[width=\textwidth]{component-diagram.png}
    \caption{Component Diagram}
    \label{fig:component-diagram}
\end{figure}

\textbf{Mikroservisler:}
\begin{enumerate}
    \item User Service
    \item Chat Service
    \item Message Service
    \item Notification Service
    \item File Service
\end{enumerate}

\textbf{Altyapı Bileşenleri:}
\begin{itemize}
    \item API Gateway
    \item WebSocket Gateway
    \item Message Queue (RabbitMQ)
    \item Cache (Redis)
    \item Databases (MongoDB/PostgreSQL)
\end{itemize}

Her servis bağımsız bir bileşen olarak modellenmiş ve kendi veritabanına sahiptir. Servisler arası iletişim REST API ve message queue üzerinden gerçekleşir.

\section{Sequence Diagram (Sıralama Diyagramı)}
Sequence diagram, servisler arasındaki mesajlaşma akışını zaman sırasına göre gösterir.

\begin{figure}[H]
    \centering
    \includegraphics[width=\textwidth]{sequence-diagram.png}
    \caption{Mesaj Gönderme Sequence Diagram}
    \label{fig:sequence-message}
\end{figure}

\begin{figure}[H]
    \centering
    \includegraphics[width=\textwidth]{sequence-diagram-file-upload.png}
    \caption{Dosya Yükleme Sequence Diagram}
    \label{fig:sequence-file}
\end{figure}

Aşağıdaki kritik senaryolar modellenmiştir:
\begin{enumerate}
    \item \textbf{Mesaj Gönderme Senaryosu:} Client $\rightarrow$ API Gateway $\rightarrow$ Message Service $\rightarrow$ Notification Service $\rightarrow$ WebSocket Gateway $\rightarrow$ Client
    \item \textbf{Dosya Yükleme Senaryosu:} Client $\rightarrow$ API Gateway $\rightarrow$ File Service $\rightarrow$ Message Service
\end{enumerate}

\section{Class Diagram (Sınıf Diyagramı)}
Class diagram, her mikroservisin iç yapısını ve sınıf ilişkilerini gösterir.

\begin{figure}[H]
    \centering
    \includegraphics[width=\textwidth]{class-diagram.png}
    \caption{Class Diagram}
    \label{fig:class-diagram}
\end{figure}

Her servis için ana sınıflar:

\textbf{User Service Sınıfları:}
\begin{itemize}
    \item UserController
    \item UserService
    \item User (Model/Entity)
\end{itemize}

\textbf{Chat Service Sınıfları:}
\begin{itemize}
    \item ChatController
    \item ChatService
    \item Chat (Model/Entity)
\end{itemize}

\textbf{Message Service Sınıfları:}
\begin{itemize}
    \item MessageController
    \item MessageService
    \item Message (Model/Entity)
\end{itemize}

\textbf{Notification Service Sınıfları:}
\begin{itemize}
    \item NotificationController
    \item NotificationService
    \item Notification (Model/Entity)
    \item WebSocketGatewayClient
\end{itemize}

\textbf{File Service Sınıfları:}
\begin{itemize}
    \item FileController
    \item FileService
    \item File (Model/Entity)
    \item StorageService
\end{itemize}

\textit{Not: Her serviste veri erişimi için ayrı bir Repository katmanı yerine, Model sınıfları direkt olarak Service katmanından kullanılmaktadır. Bu yaklaşım, basitlik ve performans açısından tercih edilmiştir.}

\section{Deployment Diagram (Dağıtım Diyagramı)}
Deployment diagram, sistemin fiziksel dağıtımını ve altyapı bileşenlerini gösterir.

\begin{figure}[H]
    \centering
    \includegraphics[width=\textwidth]{deployment-diagram.png}
    \caption{Deployment Diagram}
    \label{fig:deployment-diagram}
\end{figure}

Sistem aşağıdaki düğümlerde (nodes) çalışır:

\textbf{Container Düğümleri:}
\begin{itemize}
    \item API Gateway Container
    \item WebSocket Gateway Container
    \item 5 Mikroservis Container'ları
\end{itemize}

\textbf{Altyapı Düğümleri:}
\begin{itemize}
    \item Message Queue Node (RabbitMQ)
    \item Cache Node (Redis)
    \item Database Nodes (MongoDB/PostgreSQL)
    \item Storage Node (File Storage)
\end{itemize}

Her düğüm, containerization (Docker) ile yönetilir ve orchestration (Docker Compose) ile koordine edilir.

\section{Tasarım Prensipleri}

\subsection{Mikroservis Prensipleri}
\begin{itemize}
    \item \textbf{Bağımsız Dağıtılabilirlik:} Her servis bağımsız olarak dağıtılabilir.
    \item \textbf{Tek Sorumluluk:} Her servis belirli bir iş alanından sorumludur.
    \item \textbf{Veritabanı Ayrımı:} Her servis kendi veritabanına sahiptir.
    \item \textbf{API Tabanlı İletişim:} Servisler REST API ve message queue ile iletişim kurar.
\end{itemize}

\subsection{Tasarım Desenleri}
\begin{itemize}
    \item \textbf{API Gateway Pattern:} Tüm dış istekler API Gateway üzerinden yönlendirilir.
    \item \textbf{Event-Driven Architecture:} Asenkron işlemler için message queue kullanılır.
    \item \textbf{Circuit Breaker:} Servis hatalarına karşı dayanıklılık sağlar.
\end{itemize}
